\begin{remark}
\label{remark-adjunction}
We will use frequently the standard adjunction between $\Hom$ and tensor 
product, in the form of the natural isomorphism of contravariant functors
\begin{equation}
\label{equation-adjunction}
C(\bullet_1 \otimes_R \bullet_2) \cong \Hom_R(\bullet_1, C(\bullet_2)): 
\text{Mod}_R \times \text{Mod}_R \to \text{Mod}_R
\end{equation}
taking $f: M_1 \otimes_R M_2 \to \mathbf{Q}/\mathbf{Z}$ to the map $m_1 \mapsto 
(m_2 \mapsto f(m_1 \otimes m_2))$. See
Algebra, Lemma \ref{algebra-lemma-hom-from-tensor-product-variant}.
A corollary of this observation is that if
$$
\xymatrix@C=9pc{
C(M) \ar@<1ex>[r] \ar@<-1ex>[r] & C(N) \ar[r] & C(P)
}
$$
is a split coequalizer diagram in $\text{Mod}_R$, then so is
$$
\xymatrix@C=9pc{
C(M \otimes_R Q) \ar@<1ex>[r] \ar@<-1ex>[r] & C(N \otimes_R Q) \ar[r] & C(P 
\otimes_R Q)
}
$$
for any $Q \in \text{Mod}_R$.
\end{remark}